\documentclass[11pt]{article}
\usepackage[margin=1in]{geometry}
\usepackage{amsmath,amssymb,amsthm}
\usepackage{xurl}
\usepackage{hyperref}
\sloppy
\let\oldus\_ \renewcommand{\_}{\oldus\allowbreak}
\newtheorem{theorem}{Theorem}
\newtheorem{lemma}[theorem]{Lemma}
\newtheorem{corollary}[theorem]{Corollary}
\theoremstyle{definition}
\newtheorem{definition}[theorem]{Definition}
\newcommand{\Ext}{\operatorname{Ext}}
\newcommand{\depth}{\operatorname{depth}}

\title{Proof of conjecture 00000003200:\\ maximal regular sequences have length equal to the Ext-depth (Rees)}
\author{Jackmeson1}
\date{2026-10-05}

\begin{document}
\maketitle

\section{The conjecture and the reading}

\textbf{Statement (English, verbatim).} Definition: Dimension theory: Krull dimension and Noether
normalization. Conjecture: The linear-algebraic computation of dimension: the maximal length of
regular sequences equals the depth of the homogeneous ideal, with the evaluation of depth the Ext
vanishing layer. (Noether normalization Ext evaluation)

The Chinese text says the same thing: the maximal length of regular sequences equals the depth of the
homogeneous ideal, and depth is evaluated as the Ext vanishing layer.

\textbf{Reading.} The claim is Rees's theorem: the maximal length of a regular sequence equals the
depth, where the depth is the first degree in which Ext does not vanish. We prove it in the general
form below and then specialize it to a standard graded polynomial ring and a homogeneous ideal. The
phrases ``Krull dimension'', ``Noether normalization'' and ``linear-algebraic computation'' in the
Definition line play no role in the claim. We do not claim $\depth=\dim$ (the Cohen--Macaulay
property), since the text does not state it.

\textbf{Setting and conventions.} $R$ is a commutative Noetherian ring, $J\subseteq R$ an ideal and
$M$ a finitely generated $R$-module with $JM\neq M$ (in Lean, $J\cdot\top<\top$ in the submodule
lattice of $M$; $M$ is an object of \texttt{ModuleCat} in the universe of $R$).
\begin{itemize}
\item A sequence $x_1,\dots,x_n\in R$ is \emph{$M$-regular} if each $x_i$ is a nonzerodivisor on
  $M/(x_1,\dots,x_{i-1})M$ and $M\neq(x_1,\dots,x_n)M$ (Mathlib's
  \texttt{RingTheory.Sequence.IsRegular}). The empty sequence is $M$-regular iff $M\neq0$.
\item It is a regular sequence \emph{in $J$} if all $x_i\in J$ (\texttt{IsRegIn}). It is
  \emph{maximal} if no $r\in J$ makes $x_1,\dots,x_n,r$ an $M$-regular sequence
  (\texttt{IsMaximalRegIn}).
\item $\Ext^i_R(R/J,M)$ is Mathlib's \texttt{CategoryTheory.Abelian.Ext} in \texttt{ModuleCat R};
  ``$\neq0$'' means not \texttt{Subsingleton}.
\item $\texttt{extDepth}\,J\,M=\min\{i\in\mathbb N:\Ext^i_R(R/J,M)\neq0\}\in\mathbb N\cup\{\infty\}$
  (infimum in $\mathbb N_\infty$, $=\infty$ if the set is empty). This is the $J$-depth
  $\depth_J(M)$.
\item $\texttt{regDepth}\,J\,M=\sup\{\text{length of an $M$-regular sequence in $J$}\}\in
  \mathbb N\cup\{\infty\}$; $\texttt{regLengths}\,J\,M\subseteq\mathbb N$ is the set of these lengths.
\end{itemize}

\textbf{Sources.} Wikipedia, \emph{Depth (ring theory)}: ``By a theorem of David Rees, the depth can
also be characterized using the notion of a regular sequence.'' The decisive equivalence is proved in
Mathlib, file \url{Mathlib/RingTheory/Depth/Rees.lean} (by Nailin Guan): ``In this file we prove the
Rees theorem for depth, which relates the vanishing of certain `Ext` groups and the length of a
maximal regular sequence in a certain ideal.''

\section{The general theorem}

\begin{lemma}[Rees, in Mathlib; \texttt{C3200.rees\_iff}]\label{lem:rees}
For every $n\in\mathbb N$: $\Ext^i_R(R/J,M)=0$ for all $i<n$ iff $J$ contains an $M$-regular
sequence of length $n$.
\end{lemma}
\begin{proof}
This is Mathlib's \texttt{ModuleCat.exists\_isRegular\_of\_exists\_subsingleton\_ext} and
\texttt{ModuleCat.subsingleton\_ext\_of\_exists\_isRegular} applied with $N=R/J$, whose support is
$V(J)$ since $\operatorname{Ann}(R/J)=J$ (\texttt{support\_quot}).
\end{proof}

\begin{lemma}[\texttt{C3200.length\_le\_extDepth}]\label{lem:le}
If $x_1,\dots,x_n$ is $M$-regular in $J$, then $n\le\texttt{extDepth}\,J\,M$.
\end{lemma}
\begin{proof}
By Lemma~\ref{lem:rees}, $\Ext^i(R/J,M)=0$ for $i<n$, so every nonvanishing degree is $\ge n$.
\end{proof}

\begin{lemma}[Extension; \texttt{C3200.exists\_extend}]\label{lem:ext}
If $\mathbf x=x_1,\dots,x_n$ is $M$-regular in $J$ and $\Ext^i_R(R/J,M)=0$ for all $i\le n$, then
there is $r\in J$ such that $x_1,\dots,x_n,r$ is $M$-regular.
\end{lemma}
\begin{proof}
Induction on $n$, for all $M$ at once. For $n=0$, Lemma~\ref{lem:rees} with $n=1$ gives a regular
sequence $[r]$ in $J$. For $n\ge1$, put $a=x_1$ and $M'=M/aM$. Then $x_2,\dots,x_n$ is $M'$-regular,
and $JM'\neq M'$ (a lemma copied from a private lemma of Mathlib's \texttt{Rees.lean},
\texttt{smul\_top\_quot\_ne\_top}). The short exact sequence $0\to M\xrightarrow{a}M\to M'\to0$
gives an exact segment
$\Ext^i(R/J,M)\to\Ext^i(R/J,M')\to\Ext^{i+1}(R/J,M)$ (Mathlib's
\texttt{Ext.covariant\_sequence\_exact}$_3$\texttt{'}). For $i\le n-1$ both ends vanish, so
$\Ext^i(R/J,M')=0$. The induction hypothesis extends $x_2,\dots,x_n$ by some $r\in J$ on $M'$, and
then $a,x_2,\dots,x_n,r$ is $M$-regular.
\end{proof}

\begin{theorem}[\texttt{C3200.length\_eq\_extDepth\_of\_maximal}]\label{thm:max}
Every maximal $M$-regular sequence in $J$ has length $\texttt{extDepth}\,J\,M$.
\end{theorem}
\begin{proof}
Let $\mathbf x$ be maximal of length $n$. Lemma~\ref{lem:le} gives $n\le\texttt{extDepth}$. If
$\Ext^n(R/J,M)=0$, then together with Lemma~\ref{lem:rees} (degrees $<n$) Lemma~\ref{lem:ext}
extends $\mathbf x$, a contradiction. So $\Ext^n\neq0$ and $\texttt{extDepth}\le n$.
\end{proof}

\begin{lemma}[\texttt{C3200.exists\_maximal}]\label{lem:exists}
A maximal $M$-regular sequence in $J$ exists.
\end{lemma}
\begin{proof}
$M\neq0$, so the empty sequence is regular. Since $R$ is Noetherian
(\texttt{set\_has\_maximal\_iff\_noetherian}), choose a regular sequence $\mathbf x$ in $J$ whose ideal
$(\mathbf x)$ is maximal among the ideals of such sequences. If $\mathbf x,r$ were regular with
$r\in J$, then $(\mathbf x)=(\mathbf x,r)$, so $r\in(\mathbf x)$ acts as $0$ on
$\overline M=M/(\mathbf x)M$, while it is a nonzerodivisor on $\overline M$. Hence $\overline M=0$,
contradicting $M\neq(\mathbf x)M$.
\end{proof}

\begin{theorem}[Rees; \texttt{C3200.rees\_depth}]\label{thm:main}
There is $n\in\mathbb N$ with
(a) $\texttt{extDepth}\,J\,M=n$;
(b) $\texttt{regDepth}\,J\,M=n$;
(c) $n$ is the greatest element of $\texttt{regLengths}\,J\,M$ (\texttt{IsGreatest});
(d) every maximal $M$-regular sequence in $J$ has length $n$.
\end{theorem}
\begin{proof}
Take $\mathbf x$ from Lemma~\ref{lem:exists} and $n$ its length. Then (a) and (d) are
Theorem~\ref{thm:max}. Lemma~\ref{lem:le} bounds every length by $n$, and $\mathbf x$ attains $n$,
which gives (b) and (c).
\end{proof}

\section{The graded readings}

Let $k$ be a field and $S=k[x_1,\dots,x_m]$ (\texttt{MvPolynomial (Fin m) k}) with the standard
grading by total degree (Mathlib's \texttt{MvPolynomial.gradedAlgebra} on
\texttt{homogeneousSubmodule}). ``Homogeneous'' is Mathlib's \texttt{Ideal.IsHomogeneous} for this
grading. The irrelevant ideal $\mathfrak m=\texttt{irr}\,k\,m$ is Mathlib's
\texttt{HomogeneousIdeal.irrelevant}. We prove $f\in\mathfrak m\iff$ the constant coefficient of $f$
is $0$ (\texttt{mem\_irr\_iff}), so $\mathfrak m=(x_1,\dots,x_m)$ is the kernel of $S\to k$
(\texttt{irr\_eq\_ker}). It is maximal (\texttt{irr\_isMaximal}), and $S/\mathfrak m=k$ is the
residue field.

\begin{lemma}[\texttt{C3200.le\_irr}]
A proper homogeneous ideal $I$ of $S$ is contained in $\mathfrak m$.
\end{lemma}
\begin{proof}
If $f\in I$, its degree-$0$ component, the constant $c$, lies in $I$. If $c\neq0$ it is a unit, so
$I=S$.
\end{proof}

\begin{corollary}[Reading R1, grade; \texttt{C3200.grade\_reading}]
For a proper ideal $I$ of $S$ (in particular a proper homogeneous one), Theorem~\ref{thm:main} holds
with $J=I$ and $M=S$. So the maximal length of an $S$-regular sequence in $I$ equals
$\min\{i:\Ext^i_S(S/I,S)\neq0\}=\operatorname{grade}I$, and every maximal one has this length.
\end{corollary}

\begin{corollary}[Reading R2, depth of $S/I$; \texttt{C3200.quotient\_reading}]
For a proper homogeneous ideal $I$ of $S$, Theorem~\ref{thm:main} holds with $J=\mathfrak m$ and
$M=S/I$. So the maximal length of an $S/I$-regular sequence in $\mathfrak m$ equals
$\min\{i:\Ext^i_S(k,S/I)\neq0\}$, which is Ext from the residue field $k=S/\mathfrak m$.
\end{corollary}
\begin{proof}
$\mathfrak m(S/I)=S/I$ would give $\mathfrak m+I=S$. But $\mathfrak m+I=\mathfrak m\neq S$ by the
lemma.
\end{proof}

\textbf{Scope.} The regular sequences are arbitrary elements of $J$. Homogeneous regular sequences are
not treated separately, and we do not instantiate the depth of $I$ viewed as a module (the general
Theorem~\ref{thm:main} applies to any finitely generated $M$ with $JM\neq M$, but we did not verify
that hypothesis for $M=I$).

\section{Lean formalization and axiom audit}

File \texttt{lean/Conjecture3200/Basic.lean} (257 lines, \texttt{import Mathlib} only, namespace
\texttt{C3200}). Main theorems: \texttt{rees\_depth}, \texttt{grade\_reading},
\texttt{quotient\_reading}, \texttt{irr\_isMaximal}. The main statement reads:
\begin{verbatim}
theorem rees_depth (J : Ideal R) (M : ModuleCat.{u} R) [Module.Finite R M]
    (hJ : J * (top : Submodule R M) < top) :
    exists n : Nat, extDepth J M = n /\ regDepth J M = n /\
      IsGreatest (regLengths J M) n /\
      forall rs, IsMaximalRegIn J M rs -> rs.length = n
\end{verbatim}
(ASCII rendering. In Lean, \texttt{*} is the scalar action $\bullet$ and \texttt{top} is $\top$, with
\texttt{[CommRing R] [IsNoetherianRing R]}.) \texttt{lake build} succeeds, and
\texttt{\#print axioms} reports only \texttt{propext}, \texttt{Classical.choice} and
\texttt{Quot.sound} for all four theorems. There is no \texttt{sorry} and no \texttt{native\_decide}.
The key equivalence and the private lemma \texttt{smul\_top\_quot\_ne\_top} are reused from Mathlib's
\texttt{Rees.lean} with credit.

\end{document}
